\documentclass[12pt,a4paper,oneside]{article}
\usepackage[brazil]{babel}
\usepackage[utf8]{inputenc}
\usepackage{olymp}
\usepackage{color}
\usepackage[dvipsnames]{xcolor}
\usepackage{listings}

\setcounter{problem}{3}

\begin{document}

\begin{problem}{Ordenar}{ordenar}{2}
 
O seguinte programa foi escrito para ler três números inteiros e escrevê-los em ordem crescente. Após ler os valores, uma função \texttt{ordenar} é chamada e em seguida os valores são impressos.

\lstset{language=C++,
                basicstyle=\ttfamily,
                keywordstyle=\color{Mulberry}\ttfamily,
                stringstyle=\color{Maroon}\ttfamily,
                commentstyle=\color{YellowGreen}\ttfamily,
                morecomment=[l][\color{magenta}]{\#},
                basewidth = {.48em},
                showstringspaces=false
}
{\small
\begin{lstlisting}
#include <iostream>
using namespace std;

int main()
{
    int a, b, c;
    cin >> a >> b >> c;
    ordenar(a, b, c);
    cout << a << " " << b << " " << c << endl;
    return 0;
}
\end{lstlisting}
}

O problema é que não existe a função \texttt{ordenar} usada neste programa. Sua tarefa é completar o programa criando esta função. A função deve seguir o protótipo abaixo. Ela recebe os três números e deve retorná-los por referência em ordem crescente, de modo que o programa funcione.

\texttt{\textcolor{Mulberry}{void} ordenar(\textcolor{Mulberry}{int \&}, \textcolor{Mulberry}{int \&}, \textcolor{Mulberry}{int \&});}

%\Notes
%
%Observações, se tiver.

\InputFile

A entrada contém três valores inteiros, $A$, $B$ e $C$. Restrição: $1 \leq A, B, C \leq 10^6$.


\OutputFile

Seu programa deve gerar apenas uma linha de saída, contendo os três valores em ordem crescente.

\Notes
Nesta questão, seu programa deve usar a função \texttt{main} dada acima exatamente como está, não a modifique! Sua tarefa é apenas criar a função \texttt{ordenar}.

\Example

\begin{example}
\exmp{
20 30 40
}{
20 30 40
}%
\end{example}

\begin{example}
\exmp{
84 56 89
}{
56 84 89
}%
\end{example}

\begin{example}
\exmp{
9 8 7
}{
7 8 9
}%
\end{example}

\begin{example}
\exmp{
50 9  50
}{
9 50 50
}%
\end{example}


%Comentários sobre o exemplo, se necessário.

\end{problem}

\end{document}
